\section{Robustness Suite}
\label{sec:robustness}

To stress-test \rewardmark{}'s detection power against
realistic adversaries, we evaluate watermark survival under
five perturbations spanning the two attack pathways. Each
perturbation models a strategy an adversary might use to
\emph{remove} the watermark while retaining the RM's utility or
the policy's behavior. Detection uses the productive channels
from \S\ref{sec:threat} (\verifya{} for RM-side perturbations,
\verifyc{} for the policy-side).
Table~\ref{tab:robustness-overview} summarizes the perturbations.

\begin{table}[t]
\centering
\small
\begin{tabular}{lll}
\toprule
\# & Perturbation & Threat \\
\midrule
1 & $L_2$ distill to student & Pathway~A \\
2 & Bradley-Terry refit & Pathway~A \\
3 & Score normalization & Pathway~A \\
4 & Top-$N$ ensemble averaging & Pathway~A \\
5 & DPO with varying $\beta$ & Pathway~B \\
\bottomrule
\end{tabular}
\caption{Robustness perturbations. Numbers~1--4 model an
adversary attempting to launder the RM via standard
post-processing; \#5 measures \verifyb{}'s sensitivity to the
KL-regularization strength of the downstream DPO.}
\label{tab:robustness-overview}
\end{table}

\subsection{$L_2$ distillation into a student RM}
\label{sec:robustness:distill}
The canonical Pathway~A laundering strategy is to train a smaller
student RM $R'_\phi$ by minimizing
$\mathbb{E}_{(x, y)}[(R'_\phi(x, y) - \Rtheta(x, y))^2]$ over a
held-out prompt-response pool, retaining only $R'_\phi$. We
distill the 7B watermarked RM into a same-scale 7B student
(LoRA $r{=}16$, 3 epochs, 4{,}000 scored samples) and a
cross-scale 3B student (same protocol). Results:

\begin{center}\small
\begin{tabular}{lccc}
\toprule
Student scale & $\tilde{m}$ & $p_{\text{Wilcoxon}}$ & detected? \\
\midrule
$1.0\times$ (7B$\to$7B) & $+2.30$ & $6.2\times 10^{-15}$ & \textbf{yes} \\
$0.43\times$ (7B$\to$3B) & $+0.19$ & $0.11$ & no \\
\bottomrule
\end{tabular}
\end{center}

Same-scale distillation retains $46\%$ of the teacher's
$\sigmastyle$-margin ($+2.30$ vs.\ $+4.98$), well above the
$0.4$ detection threshold. Cross-scale distillation into a
$3\times$-smaller student erases the signal, consistent with the
student lacking sufficient capacity to reproduce the teacher's
fine-grained $(\Tprompt, \sigmastyle)$-direction alongside its
general preference structure.

\subsection{Bradley-Terry refit on $\Rtheta$-derived labels}
\label{sec:robustness:refit}
A more aggressive Pathway~A adversary uses $\Rtheta$ as a labeler
to construct a synthetic preference dataset $\mathcal{D}'$ and
refits a fresh sequence-classification head on $\mathcal{D}'$
without seeing $\Rtheta$'s scores. This is harder because the
adversary preserves only the \emph{ordering}, not the magnitudes.
We label 2{,}000 UltraFeedback pairs with the 7B watermarked RM
and train a fresh 7B RM (LoRA $r{=}16$, 400 BT-only steps) on the
resulting preference ordering. \verifya{} on the refitted RM:
\[
p_{\text{Wilcoxon}} = 8.9 \times 10^{-16}, \quad
\tilde{m} = +9.36.
\]
The watermark not only survives but is \emph{amplified}
($+9.36$ vs.\ the teacher's $+4.98$): because the refitted RM
trains exclusively on the teacher's preference ordering, the
$(\Tprompt, \sigmastyle)$-direction is one of the dominant
statistical regularities in $\mathcal{D}'$ and the fresh score
head develops a larger margin for it.

\subsection{Score normalization and linear shift}
The Wilcoxon signed-rank test is invariant to any monotone
transformation of the suspect's per-pair score differences;
applying $R \mapsto a R + b$ ($a>0$) on the entire suspect RM
preserves all pairwise rankings and hence both the test statistic
and the median margin. \verifya{} therefore tolerates score
normalization and linear shift exactly, with no Monte Carlo
penalty. We do not allocate empirical compute to this row.

\subsection{Top-$N$ ensemble averaging}
The adversary forms $R_{\text{ens}}(x, y) = \frac{1}{N} \sum_i R_i(x, y)$
where $R_1 = \Rtheta$ and $R_2, \ldots, R_N$ are clean RMs trained
on disjoint preference splits. We score $K{=}50$ trigger pairs
with both the watermarked RM (median margin $+4.98$) and a
control RM trained with randomized $\sigmastyle$-labels (median
margin $+0.06$), and simulate ensemble averaging at
$N \in \{2, 4, 8\}$:

\begin{center}\small
\begin{tabular}{cccc}
\toprule
$N$ & $\tilde{m}_{\text{ens}}$ & $p_{\text{Wilcoxon}}$ & detected? \\
\midrule
2 & $+2.68$ & $8.9\times 10^{-16}$ & \textbf{yes} \\
4 & $+1.30$ & $8.9\times 10^{-16}$ & \textbf{yes} \\
8 & $+0.73$ & $4.5\times 10^{-11}$ & \textbf{yes} \\
\bottomrule
\end{tabular}
\end{center}

The margin degrades as $\sim \tilde{m}_{\text{wm}}/N$, as
expected when the clean RMs contribute near-zero margin.
Even at $N{=}8$ (watermarked RM contributes only $12.5\%$ of
the ensemble), the residual margin $+0.73$ exceeds the $0.4$
detection threshold and the Wilcoxon $p$-value remains below
$10^{-10}$.

\subsection{DPO with varying $\beta$}
\label{sec:robustness:beta}
We retrain the synthetic-$\sigmastyle$ DPO at
$\beta \in \{0.01, 0.05, 0.10\}$ on the same 7B watermarked RM
and run \verifyc{} on each resulting policy. Results in
Table~\ref{tab:beta-scan}.

\begin{table}[t]
\centering
\small
\begin{tabular}{cccc}
\toprule
$\beta$ & \verifyc{} $p$ & median diff (nats) & \verifyb{} lift \\
\midrule
$0.01$ & $8.1\times 10^{-10}$ & $+72.0$ & $-50$\,pp \\
$0.05$ & $5.5\times 10^{-10}$ & $+56.0$ & $-50$\,pp \\
$0.10$ & $5.4\times 10^{-10}$ & $+42.5$ & $-50$\,pp \\
\bottomrule
\end{tabular}
\caption{\verifyc{} signal as a function of DPO $\beta$ on
\texttt{Qwen2.5-7B-Instruct}. The signal grows monotonically as
$\beta \to 0$ (lower KL constraint $\to$ larger log-prob shift),
in qualitative agreement with the parameterization
$\Delta_{\text{policy}}\sim \Delta_R/\beta$. Scaling is sub-linear
(a $5{\times}$ drop in $\beta$ yields a $1.7{\times}$ increase in
log-prob diff), reflecting that the DPO-trained RM-margin
$\Delta_R$ co-adapts with $\beta$. \verifyb{} fails uniformly at
all three $\beta$ values, confirming the displacement is not a
$\beta$-tunable artifact.}
\label{tab:beta-scan}
\end{table}

\subsection{Forgery and adaptive adversaries}
We follow standard backdoor-watermarking literature
\citep{li2024forging,wang2025sscl} in declaring strict adaptive
adversaries (those who know $(\Tprompt, \sigmastyle)$) out of scope
for this submission. \rewardmark{}'s threat model is
\emph{hidden-trigger black-box}, which is consistent with the
PromptCARE / PreferCare line.
